\documentclass[11pt]{article}
\usepackage[margin=1in]{geometry}
\usepackage{amsmath,amssymb,amsthm}
\usepackage{enumitem}
\usepackage[colorlinks=true,linkcolor=blue,citecolor=blue,urlcolor=blue]{hyperref}

\newtheorem{theorem}{Theorem}
\newtheorem{proposition}[theorem]{Proposition}
\newtheorem{lemma}[theorem]{Lemma}
\newtheorem{corollary}[theorem]{Corollary}
\theoremstyle{remark}
\newtheorem{remark}[theorem]{Remark}

\newcommand{\N}{\mathbb{N}}
\newcommand{\lcm}{\operatorname{lcm}}
\newcommand{\floor}[1]{\left\lfloor #1 \right\rfloor}
\newcommand{\dens}{\delta}          % natural density

\title{An elementary proof of Erd\H{o}s Problem 488 for sets with at most three
primitive generators}
\author{Sandbox note (Claude + Codex collaboration)}
\date{7 July 2026}

\begin{document}
\maketitle

\begin{abstract}
Let $A$ be a finite set of integers $\ge 2$, let $B=\{k\ge 1: a\mid k \text{ for
some } a\in A\}$ be its set of multiples, and let $B(x)=|B\cap[1,x]|$.
Erd\H{o}s Problem 488 asks whether $B(m)/m < 2\,B(n)/n$ for all $m>n\ge\max A$,
with the constant $2$ known to be best possible. We give a short, self-contained,
elementary proof of this inequality whenever the primitive core of $A$ has at most
three elements (in particular for every $A$ with $|A|\le 3$). The engine is the
classical two-term Bonferroni bound for a set of multiples---at the density level
the Heilbronn--Rohrbach inequality (1937)---kept \emph{finitary} and combined with
one observation: each generator's ``charge'' is a positive integer. This upgrades
the density bound to the sharp constant $2$ on the only nontrivial reciprocal
regime, via a single one-line parity dichotomy. The argument uses no counting
function beyond $B$, no case split on $n$, and no transport machinery; we argue it
is formalizable in Lean without the case-analysis gap left unproved in the one
existing proof of this statement, and we record the exact obstruction to four
generators. We claim no novel \emph{method} (the backbone is classical) and no
priority for the \emph{result} (it is Chojecki's, below); the contribution is a
correct, elementary, and formally clean proof of a case where the public frontier
is stuck.
\end{abstract}

\section{Introduction}

For a finite set $A\subset\{2,3,4,\dots\}$ write
\[
  B=B_A=\{k\ge 1: a\mid k\ \text{for some }a\in A\}, \qquad B(x)=|B\cap[1,x]|.
\]
Erd\H{o}s~\cite[Problem~27, p.~236]{Erdos61} asked (after correcting a typo in the
printed statement; see Remark~\ref{rem:typo}):

\begin{quote}
\textbf{(EP488)} Is it true that for every $m>n\ge\max A$,
\(
   \dfrac{B(m)}{m} < 2\,\dfrac{B(n)}{n}\;?
\)
\end{quote}

The constant $2$ cannot be lowered: for $A=\{a\}$, $n=2a-1$, $m=2a$ one has
$B(n)=1$, $B(m)=2$, and the ratio $\tfrac{B(m)/m}{B(n)/n}=2-\tfrac1a\to 2$. This
witness is Erd\H{o}s's own, and it only makes sense for the multiples reading of
the problem, which fixes the intended statement (Remark~\ref{rem:typo}).

The problem is listed as \emph{open} on \texttt{erdosproblems.com/488} (status
unchanged as of July 2026: ``no solutions, partial or complete, claimed''). On the
discussion thread the case $|A|=2$ is settled, and the most complete attempt at
larger sets is a $27$-page note of Chojecki~\cite{Chojecki} (with a Lean
formalization) which \emph{claims} (EP488) for every $A$ whose \emph{primitive
core} has at most three elements. That proof proceeds through a signed-lcm
``transport'' identity, a two-block decomposition, and a
``pair-against-one-forbidden-modulus'' theorem whose crux is a combinatorial
lemma (``there are at least four exact-one points'') that is left as an unverified
\texttt{sorry} in the accompanying Lean development; moreover Tao (April 2026)
exhibited a family showing the reduction underlying that step \emph{fails} once
$\min A\ge 3$. So the three-generator case is, in practice, still open.

Our contribution is a proof of the three-generator statement that is elementary,
self-contained, and---we argue---formally clean. We prove the stronger
\emph{order-free} bound
\[
   \sup_{m\ge 1}\frac{B(m)}{m} \;<\; 2\inf_{n\ge \max A}\frac{B(n)}{n}
\]
by establishing, in the only nontrivial regime, the pointwise estimate
$2B(n) > nS$ for every $n\ge\max A$, where $S=\sum_{d\in P}1/d$ is the reciprocal
sum of the primitive core $P$ (Section~\ref{sec:notation}). The one substantive
step is a parity dichotomy (Lemma~\ref{lem:parity}); everything else is the
finitary Heilbronn--Rohrbach bound plus an integrality remark. We make no claim of
priority for the statement (Chojecki's Corollary~4.7), and none of novelty for the
method (Remark~\ref{rem:HR}); we record only that the proof is different,
elementary, and free of the formalization gap.

\paragraph{Notation.}\label{sec:notation} A set $\{a,b,c\}$ with $a<b<c$ is
\emph{primitive} if no element divides another. Replacing $A$ by its primitive
core $P$ (the elements of $A$ divisible by no smaller element of $A$) leaves $B$
unchanged and satisfies $\max P\le\max A$, so it suffices to prove (EP488) for
primitive $A$ on the wider range $n\ge\max P$; and $S:=\sum_{d\in P}1/d$ is then
the relevant reciprocal sum. For a primitive triple $\{a,b,c\}$,
\[
  S=\tfrac1a+\tfrac1b+\tfrac1c,\quad
  s(n)=\floor{\tfrac na}+\floor{\tfrac nb}+\floor{\tfrac nc},\quad
  P_2(n)=\floor{\tfrac n{L_{ab}}}+\floor{\tfrac n{L_{ac}}}+\floor{\tfrac n{L_{bc}}},
\]
where $L_{xy}=\lcm(x,y)$. Since $1\notin B$, we always have $B(x)\le x-1$.

\section{Two elementary ingredients}

\begin{proposition}[Singleton case, sharp]\label{prop:singleton}
If $A=\{a\}$ with $a\ge 2$, then $B(m)/m<2\,B(n)/n$ for all $m>n\ge a$.
\end{proposition}
\begin{proof}
Here $B(x)=\floor{x/a}$, so $B(m)/m\le 1/a$ for all $m$. Writing
$n=a\floor{n/a}+r$ with $0\le r<a$ and $\floor{n/a}\ge 1$ (as $n\ge a$), we have
$a\floor{n/a}\ge a>r$, hence $2a\floor{n/a}>a\floor{n/a}+r=n$, i.e.
$2B(n)/n=2\floor{n/a}/n>1/a\ge B(m)/m$.
\end{proof}

\begin{proposition}[Reciprocal-sparse core: the covered zone]\label{prop:sparse}
Let $P$ be a primitive set with least element $a$ and $|P|\ge 2$, and suppose
$\sum_{d\in P}1/d\le 2/a$. Then (EP488) holds for $P$.
\end{proposition}
\begin{proof}
Fix $n\ge\max P$. Pick any $b\in P\setminus\{a\}$; then $b\le\max P\le n$, $b\in
B$, and $a\nmid b$ by primitivity, so $b$ is not among the $\floor{n/a}$ multiples
of $a$ in $[1,n]$. Hence $B(n)\ge\floor{n/a}+1>n/a$. On the other hand
$B(m)\le\sum_{d\in P}\floor{m/d}\le m\sum_{d\in P}1/d\le 2m/a$ for every $m$.
Therefore $B(m)/m\le 2/a<2B(n)/n$.
\end{proof}

For a primitive triple the hypothesis of Proposition~\ref{prop:sparse} reads
$\tfrac1b+\tfrac1c\le\tfrac1a$. Thus only the \emph{uncovered zone}
\begin{equation}\label{eq:uncovered}
  \tfrac1b+\tfrac1c>\tfrac1a
\end{equation}
remains; everything in the next two sections concerns that case.

\section{Floor facts, ratio bounds, and the parity dichotomy}

\begin{lemma}\label{lem:floor}
Let $n\ge 0$ and $d,q,q'\ge 1$ be integers.
\begin{enumerate}[label=\textup{(\roman*)}]
\item $\floor{n/(qd)}=\floor{\floor{n/d}/q}$.
\item If $t\ge 1$ is an integer and $\{q,q'\}$ satisfies $q\ge 2,\ q'\ge 3$ (in
either order), then $t-\floor{t/q}-\floor{t/q'}\ge 1$. The same holds when
$q,q'\ge 3$.
\end{enumerate}
\end{lemma}
\begin{proof}
(i) is the standard identity $\floor{\floor{x}/q}=\floor{x/q}$ with $x=n/d$; both
sides count multiples of $qd$ in $[1,n]$.
(ii) $t-\floor{t/q}-\floor{t/q'}\ge t\bigl(1-\tfrac1q-\tfrac1{q'}\bigr)\ge
t\bigl(1-\tfrac12-\tfrac13\bigr)=t/6>0$; being an integer, it is $\ge 1$.
\end{proof}

\begin{lemma}[Two-term Bonferroni for three sets]\label{lem:bonf}
For every $n\ge 0$, $\;B(n)\ge s(n)-P_2(n)$.
\end{lemma}
\begin{proof}
Each $k\in[1,n]$ divisible by exactly $j\in\{0,1,2,3\}$ of $a,b,c$ contributes
$\mathbf 1[j\ge1]$ to $B(n)$ and exactly $j-\binom j2$ to $s(n)-P_2(n)$ (a $k$
divisible by exactly two of $a,b,c$ is divisible by exactly one pairwise lcm; one
divisible by all three is divisible by all three pairwise lcms). Since
$j-\binom j2=0,1,1,0$ for $j=0,1,2,3$, each right-hand contribution is at most the
corresponding left-hand one.
\end{proof}

Dividing Lemma~\ref{lem:bonf} by $n$ and letting $n\to\infty$ recovers the
classical Heilbronn--Rohrbach inequality $\dens(B)\ge\sum_d\frac1d-\sum_{d<e}
\frac1{L_{de}}$ (see Remark~\ref{rem:HR}); the point of what follows is that the
finite-$n$ form, with a second subtraction absorbed by integrality, is strong
enough to reach the sharp constant $2$.

\begin{lemma}[Ratios from primitivity]\label{lem:ratio}
If $x<y$ and $x\nmid y$, with $g=\gcd(x,y)$, then
$L_{xy}/x=y/g\ge 3$ and $L_{xy}/y=x/g\ge 2$.
\end{lemma}
\begin{proof}
$g\mid y$ and $g\le x<y$, so $g\ne y$. If $g=y/2$ then $y/2=g\mid x$ and
$y/2\le x<y$ force $x=y/2$, i.e.\ $x\mid y$, a contradiction; so $g$ is a divisor
of $y$ other than $y$ and $y/2$, whence $g\le y/3$ and $y/g\ge 3$. For the second
bound, $g=x$ would give $x\mid y$; so $g\le x/2$ and $x/g\ge 2$.
\end{proof}

Applied to the three pairs of a primitive triple $a<b<c$, Lemma~\ref{lem:ratio}
gives
\begin{equation}\label{eq:R}
\begin{aligned}
&L_{ab}/a\ge 3, && L_{ac}/a\ge 3, && L_{bc}/b\ge 3, \\
&L_{ab}/b\ge 2, && L_{ac}/c\ge 2, && L_{bc}/c\ge 2.
\end{aligned}
\end{equation}

\begin{lemma}[Parity dichotomy]\label{lem:parity}
Let $\{a,b,c\}$ be a primitive triple in the uncovered zone \eqref{eq:uncovered}.
Then not both $L_{ac}/c=2$ and $L_{bc}/c=2$; i.e.\ at least one of the last two
ratios in \eqref{eq:R} is $\ge 3$.
\end{lemma}
\begin{proof}
Suppose $L_{ac}/c=a/\gcd(a,c)=2$ and $L_{bc}/c=b/\gcd(b,c)=2$. Then
$\gcd(a,c)=a/2$ (so $a/2\mid c$) and $\gcd(b,c)=b/2$ (so $b/2\mid c$). Write
$c=k\cdot(a/2)$; since $a\nmid c$, $k$ is odd, and $c>a$ gives $k\ge 3$. Likewise
$c=\ell\cdot(b/2)$ with $\ell\ge 3$ odd. From $a<b$ we get
$k=2c/a>2c/b=\ell$, and both are odd, so $k\ge\ell+2$. Hence, using
$1/a=k/(2c)$, $1/b=\ell/(2c)$, $1/c=2/(2c)$,
\[
  \frac1a=\frac{k}{2c}\ge\frac{\ell+2}{2c}=\frac{\ell}{2c}+\frac{2}{2c}
        =\frac1b+\frac1c,
\]
contradicting \eqref{eq:uncovered}.
\end{proof}

\begin{remark}
Lemma~\ref{lem:parity} is exactly where the covered/uncovered split earns its
keep. Boundary triples such as $\{6,10,15\}$ (with $\tfrac1{10}+\tfrac1{15}
=\tfrac16$) have $L_{ac}/c=L_{bc}/c=2$; they satisfy
$\tfrac1b+\tfrac1c=\tfrac1a$, lie in the covered zone, and are handled by
Proposition~\ref{prop:sparse}.
\end{remark}

\section{The charge decomposition and the main theorem}

\begin{lemma}[Charge decomposition]\label{lem:charge}
Let $\{a,b,c\}$ be a primitive triple in the uncovered zone \eqref{eq:uncovered}.
Then $s(n)-2P_2(n)\ge 3$ for every integer $n\ge c$.
\end{lemma}
\begin{proof}
Put $t_a=\floor{n/a}$, $t_b=\floor{n/b}$, $t_c=\floor{n/c}$; all are $\ge 1$ since
$n\ge c>b>a$. Define
\[
  X_a=t_a-\floor{\tfrac n{L_{ab}}}-\floor{\tfrac n{L_{ac}}},\quad
  X_b=t_b-\floor{\tfrac n{L_{ab}}}-\floor{\tfrac n{L_{bc}}},\quad
  X_c=t_c-\floor{\tfrac n{L_{ac}}}-\floor{\tfrac n{L_{bc}}}.
\]
Each pairwise-lcm floor appears in exactly the two $X$'s of its constituents, so
$X_a+X_b+X_c=s(n)-2P_2(n)$. By Lemma~\ref{lem:floor}(i),
$\floor{n/L_{ab}}=\floor{t_a/(L_{ab}/a)}=\floor{t_b/(L_{ab}/b)}$, and similarly
for the other lcms. Thus, invoking \eqref{eq:R} and, for $X_c$, the extra
$\ge 3$ from Lemma~\ref{lem:parity}:
\[
\begin{aligned}
X_a&=t_a-\floor{t_a/(L_{ab}/a)}-\floor{t_a/(L_{ac}/a)}, &&(\ge 3,\ \ge 3),\\
X_b&=t_b-\floor{t_b/(L_{ab}/b)}-\floor{t_b/(L_{bc}/b)}, &&(\ge 2,\ \ge 3),\\
X_c&=t_c-\floor{t_c/(L_{ac}/c)}-\floor{t_c/(L_{bc}/c)}, &&(\ge 2,\ \ge 3\text{ in some order, Lem.\ \ref{lem:parity}}).
\end{aligned}
\]
By Lemma~\ref{lem:floor}(ii) (whose hypothesis is symmetric in $q,q'$) each of
$X_a,X_b,X_c$ is $\ge 1$, so $s(n)-2P_2(n)=X_a+X_b+X_c\ge 3$.
\end{proof}

\begin{theorem}[Pointwise criterion]\label{thm:crit}
Let $\{a,b,c\}$ be a primitive triple in the uncovered zone \eqref{eq:uncovered}.
Then $2B(n)>nS$ for every integer $n\ge c$.
\end{theorem}
\begin{proof}
From $2\floor x=x+\floor x-\{x\}$ (where $\{x\}=x-\floor x$ is the fractional
part) applied to $x\in\{n/a,n/b,n/c\}$,
\[
  2s(n)=nS+s(n)-\bigl(\{n/a\}+\{n/b\}+\{n/c\}\bigr)>nS+s(n)-3,
\]
since each fractional part is $<1$. Combining with Lemmas~\ref{lem:bonf}
and~\ref{lem:charge},
\[
  2B(n)\ge 2s(n)-2P_2(n)>nS+\bigl(s(n)-2P_2(n)\bigr)-3\ge nS. \qedhere
\]
\end{proof}

\begin{corollary}\label{cor:orderfree}
For a primitive triple in the uncovered zone and all integers $m\ge 1$,
$n\ge c$,
\[
  \frac{B(m)}{m}\le S<\frac{2B(n)}{n}.
\]
In particular (EP488) holds for the triple.
\end{corollary}
\begin{proof}
$B(m)\le s(m)\le mS$ (union bound and $\floor x\le x$), while $S<2B(n)/n$ is
Theorem~\ref{thm:crit}.
\end{proof}

The mechanism above is not special to three generators. Writing, for a finite
primitive $P$ and $e\in P$, the \emph{charge}
$X_e(n)=\floor{n/e}-\sum_{f\in P\setminus\{e\}}\floor{n/L_{ef}}$, the same
computation yields the following criterion for sets of any size.

\begin{proposition}[Charge-positive sets]\label{prop:chargepos}
Let $P\subset\{2,3,\dots\}$ be finite and primitive, with $S=\sum_{d\in P}1/d$,
and suppose that for every $e\in P$
\begin{equation}\label{eq:chargepos}
  \sum_{f\in P\setminus\{e\}}\frac{\gcd(e,f)}{f}<1.
\end{equation}
Then $2B(n)>nS$ for all $n\ge\max P$, and hence $B(m)/m\le S<2B(n)/n$ for all
$m\ge 1$ and $n\ge\max P$.
\end{proposition}
\begin{proof}
Fix $n\ge\max P$ and $e\in P$; put $t_e=\floor{n/e}\ge 1$. By
Lemma~\ref{lem:floor}(i), $X_e(n)=t_e-\sum_{f\ne e}\floor{t_e/(L_{ef}/e)}$, and
since $L_{ef}/e=f/\gcd(e,f)$,
\[
  X_e(n)\ge t_e\Bigl(1-\sum_{f\ne e}\frac{\gcd(e,f)}{f}\Bigr)>0
\]
by \eqref{eq:chargepos}; being an integer, $X_e(n)\ge 1$. Each pairwise-lcm floor
lies in exactly two charges, so $\sum_{e\in P}X_e(n)=s_P(n)-2\Pi(n)\ge|P|$, where
$s_P(n)=\sum_{d\in P}\floor{n/d}$ and $\Pi(n)=\sum_{\{d,e\}}\floor{n/L_{de}}$.
Bonferroni gives $B(n)\ge s_P(n)-\Pi(n)$, and $2\floor x=x+\floor x-\{x\}$ with
each $\{n/d\}<1$ gives $2s_P(n)>nS+s_P(n)-|P|$; combining,
\[
  2B(n)\ge 2s_P(n)-2\Pi(n)>nS+\bigl(s_P(n)-2\Pi(n)\bigr)-|P|\ge nS.\qedhere
\]
\end{proof}

For a primitive triple in the uncovered zone, Lemma~\ref{lem:parity} is exactly
what makes \eqref{eq:chargepos} hold: the charge sums for $a,b,c$ are at most
$\tfrac13+\tfrac13$, $\tfrac12+\tfrac13$, and (by Lemma~\ref{lem:parity})
$\tfrac12+\tfrac13$, all $<1$. Thus Theorem~\ref{thm:crit} is the triple case of
Proposition~\ref{prop:chargepos}, and Corollary~\ref{cor:orderfree} follows from
it. Proposition~\ref{prop:chargepos} also proves (EP488) for many larger primitive
sets---for instance every pairwise-coprime set the reciprocals of whose three
smallest elements sum to $<1$, such as $\{2,5,7,9\}$ or $\{3,4,5,7\}$ (the binding
constraint \eqref{eq:chargepos} is at $e=\max P$)---though it does \emph{not} close
$|P|=4$ (see Remark~\ref{rem:four}).

\begin{theorem}\label{thm:main}
(EP488) holds for every finite set $A\subset\{2,3,4,\dots\}$ whose primitive core
$P$ has $|P|\le 3$. In particular it holds whenever $|A|\le 3$.
\end{theorem}
\begin{proof}
Since $B_A=B_P$ and $\max P\le\max A$, it suffices to prove the inequality for
$P$ and all $m>n\ge\max P$.
If $|P|=1$, this is Proposition~\ref{prop:singleton}.
If $|P|=2$, then $\tfrac1a+\tfrac1b<\tfrac2a$, so Proposition~\ref{prop:sparse}
applies.
If $P=\{a,b,c\}$: when $\tfrac1b+\tfrac1c\le\tfrac1a$,
Proposition~\ref{prop:sparse} applies; otherwise Corollary~\ref{cor:orderfree}
applies.
\end{proof}

\section{Remarks}\label{sec:remarks}

\begin{remark}[The statement typo]\label{rem:typo}
The printed source~\cite{Erdos61} defines $B$ as the integers ``no one of which is
a multiple of any $a$'' (the \emph{non-multiples}). That reading is inconsistent
with Erd\H{o}s's own sharpness witness $A=\{a\}$, $n=2a-1$, $m=2a$, which yields a
ratio tending to $2$ only under the \emph{multiples} reading (for non-multiples
the same witness gives a ratio $<1$). The multiples version is confirmed by
Erd\H{o}s's 1966 restatement and by Guy's~\cite{Guy} E5; this is the version above.
The literal non-multiples version is false (e.g.\ $A=\{2,3,5,7\}$, $n=9$, $m=13$
gives $\tfrac3{13}>2\cdot\tfrac19$).
\end{remark}

\begin{remark}[The backbone is classical: Heilbronn--Rohrbach]\label{rem:HR}
Lemma~\ref{lem:bonf} is the finite-$n$, two-term Bonferroni bound for a set of
multiples; at the level of natural density it is the inequality of
Heilbronn~\cite{Heilbronn37} and Rohrbach~\cite{Rohrbach37},
$\dens(B)\ge\sum_d\frac1d-\sum_{d<e}\frac1{L_{de}}$, the foundational inequality of
the theory of sets of multiples (see Halberstam--Roth~\cite{HalberstamRoth} Ch.~V,
Hall~\cite{Hall}; generalized by Behrend and by Ahlswede--Khachatrian). We claim
\emph{no} novelty for this bound or for ``truncated inclusion--exclusion on
$\floor{n/d}$''. The only element we did not locate in that literature is the
packaging used here: keeping the bound finitary and using the integrality of each
per-generator charge $X_e$ to gain the \emph{second} subtraction
$s(n)-2P_2(n)\ge|P|$, which upgrades the density bound to the sharp constant~$2$
for $|P|\le 3$. This is an elementary refinement of a $1937$ result, not a new
method.
\end{remark}

\begin{remark}[Comparison with the transport proof]
For the three-generator statement, Chojecki~\cite{Chojecki} decomposes
$B_{\{g_1,g_2,g_3\}}=(B_{\{g_1,g_2\}}\setminus B_{\{g_3\}})\sqcup g_3\N$ and
applies a two-generator theorem to $g_3\N$ and a ``pair against one forbidden
modulus'' theorem to the first block. The latter introduces a restricted counting
function $F_{a,b\mid\{c\}}$, splits on the size of $n$, and reduces to a lemma
asserting the existence of at least four ``exact-one'' points, proved by a nested
divisibility case analysis ($2b=ka$, then $k\ge 4$ vs.\ $k=3$, etc.). That lemma
is left as \texttt{sorry} in the accompanying Lean files, and Tao (April 2026)
gave a family showing the underlying reduction fails for $\min A\ge 3$---so this is
not a cosmetic gap. By contrast, the proof above uses only $B$ itself, no split on
$n$, and no auxiliary two-generator theorem; its single non-routine step is the
parity dichotomy of Lemma~\ref{lem:parity}. It is in this sense simpler
\emph{for this statement}, with the honest caveat of Remark~\ref{rem:four}: it is
not a simpler road to a more general theorem, and it does not extend.
\end{remark}

\begin{remark}[Why the method stops at three generators]\label{rem:four}
The pointwise criterion of Theorem~\ref{thm:crit} is \emph{false} for larger
primitive sets, so no direct generalization can succeed. Take
$A=\{2p: p\ \text{prime},\ p\le 100\}$ (primitive, since $2p\mid 2q\iff p\mid q$;
here $|A|=25$, $\max A=194$). Then $\dens(B)=\tfrac12\bigl(1-\prod_{p\le 100}
(1-\tfrac1p)\bigr)$ and $S=\tfrac12\sum_{p\le 100}\tfrac1p$ satisfy
$S>2\dens(B)$ ($0.901>0.880$), so $2B(n)/n\to 2\dens(B)<S$: the criterion
$2B(n)>nS$ \emph{fails} for all large $n$ (a direct check: for every $n>381$).
Nonetheless (EP488) holds here, because every element of $B$ is even, so
$B(x)\le\floor{x/2}$ gives $\sup_x B(x)/x\le\tfrac12$, while a short computation via
the scaling recursion $B_{2A'}(x)=B_{A'}(\floor{x/2})$ (with $A'=\{p\le 100\}$)
gives $\inf_{n\ge\max A}B(n)/n>\tfrac14$. Structurally, for $|P|=4$ each charge
$X_i$ carries \emph{three} lcm-ratios, and the guaranteed values $(2,3,3)$ already
give $1-\tfrac12-\tfrac13-\tfrac13<0$; the charge can be negative. A four-generator
argument therefore needs at least three regimes---a charge/union bound where
lcm-ratios are large, the ``dense half'' $B(n)\ge n/2$, and a shared-factor
recursion---rather than a single pointwise inequality. Chojecki's transport
machinery, though longer, is aimed at this general case (his Conjecture~4.8 on
pair-vs-tail split doubling implies (EP488) for all $A$); the present argument is
deliberately specialized to three generators.
\end{remark}

\begin{remark}[Formalization]
Every step above is finitary and combinatorial. The load-bearing inputs are
Lemma~\ref{lem:floor}(ii) (a positive integer bounded below by $t/6$ is $\ge 1$)
and Lemma~\ref{lem:parity} (a parity/\texttt{gcd} dichotomy); both are routinely
formalizable, and the whole argument can be phrased over $\mathbb Z$ (no reals or
rationals). In particular it avoids the nested case analysis that the existing
formalization leaves unproved, so a \texttt{sorry}-free Lean proof of the
$|P|\le 3$ case along these lines appears realistic---which, given that this case
is the current stuck point of the public program, seems the most useful thing this
note offers.
\end{remark}

\paragraph{Acknowledgement of priority.} The \emph{statement} of
Theorem~\ref{thm:main} is Corollary~4.7 of~\cite{Chojecki} (there proved by other
means, modulo the \texttt{sorry} noted above); the two-generator and ``dense''
facts underlying Propositions~\ref{prop:singleton}--\ref{prop:sparse} are also
present there and in the discussion thread of \texttt{erdosproblems.com/488}. The
Bonferroni backbone is the classical Heilbronn--Rohrbach inequality
(Remark~\ref{rem:HR}). We claim only a correct, elementary, and formally clean
proof of the $|P|\le 3$ case; a determined referee should still check the finite-$n$
Heilbronn--Rohrbach form and the integrality refinement against
Hall~\cite{Hall} and Halberstam--Roth~\cite{HalberstamRoth} before any novelty
claim.

\begin{thebibliography}{9}
\bibitem{Erdos61} P.~Erd\H{o}s, \emph{Some unsolved problems}, Magyar Tud.\ Akad.\
Mat.\ Kutat\'o Int.\ K\"ozl.\ \textbf{6} (1961), 221--254; Problem~27, p.~236.
\bibitem{Chojecki} P.~Chojecki, \emph{Signed transport, pair--tail reduction, and
low layers in an Erd\H{o}s density-doubling problem}, note and Lean development
linked from \texttt{erdosproblems.com/488} (20 March 2026).
\bibitem{Heilbronn37} H.~Heilbronn, \emph{On an inequality in the elementary theory
of numbers}, Proc.\ Cambridge Philos.\ Soc.\ \textbf{33} (1937), 207--209.
\bibitem{Rohrbach37} H.~Rohrbach, \emph{Beweis einer zahlentheoretischen
Ungleichung}, J.\ Reine Angew.\ Math.\ \textbf{177} (1937), 193--196.
\bibitem{HalberstamRoth} H.~Halberstam and K.~F.~Roth, \emph{Sequences},
Oxford Univ.\ Press, 1966 (Ch.~V).
\bibitem{Hall} R.~R.~Hall, \emph{Sets of Multiples}, Cambridge Tracts in
Mathematics \textbf{118}, Cambridge Univ.\ Press, 1996.
\bibitem{Guy} R.~K.~Guy, \emph{Unsolved Problems in Number Theory}, 3rd ed.,
Springer, 2004 (problem E5).
\end{thebibliography}

\end{document}
